\documentclass{article}
\usepackage[T1]{fontenc}
\usepackage[utf8]{inputenc}
\usepackage[polish]{babel}
\usepackage{amsmath}
\usepackage{amssymb}

\begin{document}

\section*{Zadanie 1}
Niech $m(T)$ i $M(T)$ oznaczają, odpowiednio, minimalny i maksymalny stopień wyjściowy w turnieju $T$.
\begin{enumerate}
    \item[a.] Udowodnij, że jeśli $T$ jest $4n$-turniejem i $m(T) \geq n$ oraz $M(T) < 3n$, to $T$ jest silnie spójny.
    \item[b.] Dla każdego $n > 0$ podaj przykład $4n$-turnieju $T$, w którym $m(T) \geq n$, $M(T) \leq 3n$, i który nie jest silnie spójny.
\end{enumerate}

\vspace{0.5cm}
\noindent\rule{\textwidth}{0.5pt}
\vspace{0.5cm}

\section*{Zadanie 2}
Podaj wzór na funkcję $f$, taką że $f(t)$ to liczba kolorowań \textit{krawędziowych} grafu $K_4$ przy pomocy puli $t$ kolorów.

\vspace{0.5cm}
\noindent\rule{\textwidth}{0.5pt}
\vspace{0.5cm}

\section*{Zadanie 3}
Zdefiniujmy $S_p(n) = \sum_{k=1}^{p-1} k^n$. Udowodnij, że jeśli $p > 3$ jest liczbą pierwszą, to:
\begin{enumerate}
    \item[a.] istnieje nieskończenie wiele nieparzystych $n$, takich że $S_p(n) \equiv 0 \pmod{p^2}$;
    \item[b.] istnieje nieskończenie wiele parzystych $n$, takich że $S_p(n) \not\equiv 0 \pmod p$;
    \item[c.] istnieje nieskończenie wiele parzystych $n$, takich że $S_p(n) \equiv 0 \pmod p$.
\end{enumerate}
Wskazówka: Rozważ na początek $n = 2$.

\vspace{0.5cm}
\noindent\rule{\textwidth}{0.5pt}

\end{document}